\section{Introduction}
\label{SEC:Intro}

Accurate and efficient numerical methods for the simulation of complex
fluid flows are needed to reduce the design and development costs
of aerodynamic vehicles and aeropropulsion systems.
%
One of the fundamental challenges is prediction of the turbulent motion
within the flow and its effect on the overall motion of the fluid.
%
Scale-resolving methods such as
large-eddy simulation (LES) and
direct numerical simulation (DNS)
offer great promise in predicting turbulence.
%
However, due to the high Reynolds numbers involved
and resulting small turbulent scales,
very fine spatial and temporal resolution is necessary.

%%%%%%%%%%%%%%%%%%%%%%%%%%%%%%%%%%%%%%%%%%%%%%%%%%%%%%%%%%%%%%%%%%%%%%%%%%%%%%%%

For these simulations, the required cell count and computational cost
of traditional second-order numerical schemes quickly become prohibitive.
%
High-order numerical schemes can overcome this limitation.
However, such schemes have historically only been found in structured grid
methods with large stencils.
%
These methods have numerous drawbacks, such as:
difficulty in generating structured grids for complex configurations,
high sensitivity to grid quality,
complex boundary condition implementation due to the wide stencils, and
poor scalability for parallel computing.
%
These drawbacks severely limit the applicability of such methods to
complex aeropropulsion configurations.

%%%%%%%%%%%%%%%%%%%%%%%%%%%%%%%%%%%%%%%%%%%%%%%%%%%%%%%%%%%%%%%%%%%%%%%%%%%%%%%%

The flux reconstruction (FR) method provides a simple, efficient, and
easy-to-implement method for solving the Navier-Stokes equations on
unstructured grids and is accurate to an arbitrary order.
%
By employing the high-order FR method with LES
on unstructured grids, limitations related to structured grid generation,
grid resolution, and computational efficiency can be overcome.

%%%%%%%%%%%%%%%%%%%%%%%%%%%%%%%%%%%%%%%%%%%%%%%%%%%%%%%%%%%%%%%%%%%%%%%%%%%%%%%%

\subsection{Flux Reconstruction}
\label{SEC:Intro:FR}

The \key*{flux reconstruction}[Flux Reconstruction Method] method
%, sometimes referred to as the Correction Procedure via Reconstruction (CPR),
is a simple and efficient framework for high-order numerical schemes
of conservation laws.
%
It uses a set of polynomial basis-functions to represent the solution's
conserved variables within each computational cell or element of the
fluid domain.
This polynomial is used to compute the flux of the flow through the cell.
Corrections are applied at each cell face to ensure conservation.
%
The general FR framework can be used to express various other methods
such as discontinuous Galerkin (DG), spectral difference, and
spectral element methods.
%
However, the specific FR implementation in GFR is most comparable to
nodal DG methods.

%%%%%%%%%%%%%%%%%%%%%%%%%%%%%%%%%%%%%%%%%%%%%%%%%%%%%%%%%%%%%%%%%%%%%%%%%%%%%%%%

Theoretical development of the flux reconstruction method can be
found in \linecite{Huynh_2007,Huynh_2009,Huynh_2011,Huynh_2014}.
%
Details of the FR implementation in GFR and validation of the code are discussed in
\linecite{Spiegel_2015,Spiegel_2015b,Spiegel_2016}.
%
The interested reader is directed there for full details of the method.

%%%%%%%%%%%%%%%%%%%%%%%%%%%%%%%%%%%%%%%%%%%%%%%%%%%%%%%%%%%%%%%%%%%%%%%%%%%%%%%%

\subsubsection{1D Example}
\label{SEC:Intro:FR:1D_Example}

In order to provide the user with a basic understanding of how the
FR method works, we will consider the one-dimensional advection
equation:
\begin{align}
  \label{EQN:1D_Advection}
  u_t + f_x =& 0
&       u   =& u(x,t)
&       f   =& f(u)
\end{align}
where the subscripts $x$ and $t$ denote partial differentiation
with respect to space and time, respectively.
The flux $f$ is a (potentially nonlinear) function of the conserved
solution variable(s) stored in $u$.
%
The volume of the simulation domain is divided into an arbitrary
number of nonoverlapping cells, and a mapping function,
  $ r = \chi (x) $,
is used to transform each cell from physical space to a common
reference space\slash element.
%
With $r$ denoting the local coordinate variable that varies on the
reference element
  $ \stdelem = [-1,1] $,
the inverse function
  $ x = \chi^{-1} (r) $
similarly maps the reference element back to physical space.

%%%%%%%%%%%%%%%%%%%%%%%%%%%%%%%%%%%%%%%%%%%%%%%%%%%%%%%%%%%%%%%%%%%%%%%%%%%%%%%%

The fundamental steps in the FR method can briefly
be summarized as follows:
\newline
\newlength{\lenColA}
\newlength{\lenColB}
\setlength{\lenColA}{0.5\textwidth}
\setlength{\lenColB}{0.5\textwidth}
\begin{minipage}{\lenColA}
\begin{enumerate}
\item[1.]
Represent the solution within each element as a polynomial of degree $\PO$
defined by $\PO+1$ solution points.
The solution may be discontinuous across cell interfaces,
even in the absence of shocks.
\end{enumerate}
\end{minipage}%
\begin{minipage}{\lenColB}
  \insertimage[width=\textwidth]{Intro/Intro_Step1}
\end{minipage}%
\\
\begin{minipage}{\lenColA}
\begin{enumerate}
\item[2.]
Compute the \key[key:DiscontinuousFluxFunction]{discontinuous flux function}[=]
at the quadrature points.
For this example, assume those are the same as the solution points.
\end{enumerate}
\end{minipage}%
\begin{minipage}{\lenColB}
  \insertimage[width=\textwidth]{Intro/Intro_Step2}
\end{minipage}%
\\
\begin{minipage}{\lenColA}
\begin{enumerate}
\item[3.]
Compute a \key[key:CommonFluxValue]{common flux value}[=]
at the interfaces between cells using the solutions from the left and right
cells, $\left(u_L,u_R\right)$.
\end{enumerate}
\end{minipage}%
\begin{minipage}{\lenColB}
  \insertimage[width=\textwidth]{Intro/Intro_Step3}
\end{minipage}%
\\
\begin{minipage}{\lenColA}
\begin{enumerate}
\item[4.]
At each edge of the element, compute the difference between the
\key[key:DiscontinuousFluxFunction]{discontinuous flux function}[=]
and the interface \key[key:CommonFluxValue]{common flux value}[=].
These are the left\slash right flux jumps.
\end{enumerate}
\end{minipage}%
\begin{minipage}{\lenColB}
  \insertimage[width=\textwidth]{Intro/Intro_Step4}
\end{minipage}%
\\
\begin{minipage}{\lenColA}
\begin{enumerate}
\item[5.]
Use the left\slash right flux jumps from the previous step to scale the
\key[key:CorrectionFunction]{correction functions}[correction function] of
degree $\PO+1$; the scaled correction functions take the value of the flux
jump at the respective left\slash right cell boundary while leaving the
opposite cell boundary unchanged.
%
Combining the two scaled correction functions results in the corrected flux
function, which is a polynomial of degree $\PO+1$, takes the value of the flux
jumps at their respective boundaries, and approximates the zero function in
between. 
%
Note, however, that the correction functions generally have non-zero values at
the interior points.
\end{enumerate}
\end{minipage}%
\begin{minipage}{\lenColB}
  \insertimage[width=\textwidth]{Intro/Intro_Step5}
\end{minipage}%
\\
\begin{minipage}{\lenColA}
\begin{enumerate}
\item[6.]
Add the \key[key:CorrectionFunction]{corrected flux function}[correction function]
to  the \key[key:DiscontinuousFluxFunction]{discontinuous flux function}[=]
to create a \key[key:ContinuousFluxFunction]{continuous flux function}[=]
of degree $\PO+1$ that passes through the
\key[key:CommonFluxValue]{common flux values}[common flux value].
%
Note that the flux at the interior solution points is also adjusted
slightly to smoothly accommodate this change.
\end{enumerate}
\end{minipage}%
\begin{minipage}{\lenColB}
  \insertimage[width=\textwidth]{Intro/Intro_Step6}
\end{minipage}%
\\
\begin{enumerate}
\item[7.]
Compute the spatial
\key[key:ContinuousFluxFunctionDerivative]%&
    {derivative of the continuous flux function}%&
    [continuous flux function!derivative].
\item[8.]
Use the derivative of the continuous flux function to advance the
solution $u$ to the next time level.
\end{enumerate}

%%%%%%%%%%%%%%%%%%%%%%%%%%%%%%%%%%%%%%%%%%%%%%%%%%%%%%%%%%%%%%%%%%%%%%%%%%%%%%%%

For the interested reader, we elaborate on each of these steps
in the sections that follow.

%%%%%%%%%%%%%%%%%%%%%%%%%%%%%%%%%%%%%%%%%%%%%%%%%%%%%%%%%%%%%%%%%%%%%%%%%%%%%%%%

\subsubsection{Solution Polynomial}
\label{SEC:Intro:FR:Polynomial}

%%%%%%%%%%%%%%%%%%%%%%%%%%%%%%%%%%%%%%%%%%%%%%%%%%%%%%%%%%%%%%%%%%%%%%%%%%%%%%%%

At a given time $t$, the exact solution to equation \eref{EQN:1D_Advection} is
approximated locally within each cell as a polynomial of degree \PO.
%
In some instances we will use \p{\textsc{s}} to further clarify that this is the degree
of the solution polynomial.
%
This polynomial approximation provides a smooth variation of the solution within
each cell.
%
However, the polynomials between adjacent cells are generally discontinuous
across cell interfaces even in the absence of shocks.
 
%%%%%%%%%%%%%%%%%%%%%%%%%%%%%%%%%%%%%%%%%%%%%%%%%%%%%%%%%%%%%%%%%%%%%%%%%%%%%%%%

One way to represent the degree $\PO$ solution polynomial is using a
\key*{modal form}[polynomial!modal form]
\begin{align}
  \label{EQN:Modal_Form}
       u (r) =& \sum_{n=1}^{N_p} \uModal[n] (t) \psi_n (r)
& \psi_n (r) =& r^{n-1}
\end{align}
where there are $N_p=\PO+1$ modes and $\uModal[n]$ for $n=1,2,...,N_p$
represents the amplitude coefficients of the constant, linear, quadratic,
etc. polynomial modes.
%
\Fref{FIG:Polynomial_Form_Modal} illustrates the modal decomposition
of a third-order polynomial, $\PO[3]$. Note that the various polynomial modes
are symmetric\slash antisymmetric across the cell.
%
Each mode, other than the constant mode, has the following characteristics:
(1) its integral across the cell is zero, meaning it does not contribute
to the average value within the cell,
(2) the value in the center of the cell is zero, and
(3) the value at the edge of the cell is directly related to the modal
coefficient.
%
%All modes other than the constant mode
%are zero at the center of the cell, and
%the value at the edge of the cell is directly related to the modal coefficient.
%
One can clearly visualize the individual modes, but visualizing the
solution at a point requires adding the contributions of all of the modes.

%%%%%%%%%%%%%%%%%%%%%%%%%%%%%%%%%%%%%%%%%%%%%%%%%%%%%%%%%%%%%%%%%%%%%%%%%%%%%%%%

\begin{figure}
  \begin{centering}
    \begin{subfigure}{3.0in}
      \insertimage[width=\textwidth]{Poly_Form/Poly_Form_Modal}
      \caption{Modal Representation.}
      \label{FIG:Polynomial_Form_Modal}
    \end{subfigure}%
    \begin{subfigure}{3.0in}
      \insertimage[width=\textwidth]{Poly_Form/Poly_Form_Nodal}
      \caption{Nodal Representation.}
      \label{FIG:Polynomial_Form_Nodal}
    \end{subfigure}
    \caption{Modal and nodal representation of a $\PO[3]$ polynomial.}
    \label{FIG:Polynomial_Form}
  \end{centering}
\end{figure}

%%%%%%%%%%%%%%%%%%%%%%%%%%%%%%%%%%%%%%%%%%%%%%%%%%%%%%%%%%%%%%%%%%%%%%%%%%%%%%%%

Another way to represent the solution is using a
\key*{nodal form}[polynomial!nodal form]
where the degree $\PO$ polynomial is defined by $N_p=\PO+1$ solution points.
\begin{align}
  \label{EQN:Nodal_Form}
       u (r) =& \sum_{n=1}^{N_p} \uNodal[n] (t) l_n (r)
&    l_n (r) =& \prod_{\substack{ i=1 \\ i \neq n }}^{N_p}
                \frac{r-\xi_i}{\xi_n-\xi_i}
\end{align}
The coordinate values of these solution points on \stdelem{} are given by
$\xi_n$ for $n=1,2,...,N_p$.
%
For GFR, Gauss points are recommended.
%
These are symmetrically, but not uniformly, distributed across the
interior of the cell and do not lie on the edge of the cell.
%
The \key*{Lagrange polynomial}[polynomial!Lagrange] $l_n$ is of degree $\PO$
and smoothly interpolates between the $\PO+1$ nodal points, taking the value
1 at $\xi_n$ and zero at all other points $\xi_{i \neq n}$.
The solution value at each nodal point is given by $u_n=u(\xi_n)$
and is used to scale the amplitude of the corresponding Lagrange
polynomial $l_n$.
%
\Fref{FIG:Polynomial_Form_Nodal} illustrates the nodal decomposition
of the same $\PO[3]$ polynomial shown previously. One can quickly identify
the solution value at any nodal point, but visualizing the solution
at any other location requires adding all the contributions together.
In this form, one also cannot easily identify the various modes.

%%%%%%%%%%%%%%%%%%%%%%%%%%%%%%%%%%%%%%%%%%%%%%%%%%%%%%%%%%%%%%%%%%%%%%%%%%%%%%%%

The \key{modal}[polynomial!modal form]
and \key{nodal}[polynomial!nodal form] forms each have their advantages.
For example, applying a filter to remove high-frequency noise
is most readily done using the modal form.
However, setting boundary conditions at specific locations can
be easier in nodal form.
GFR's internal representation uses the nodal form, and when
necessary switches to modal form through the use of a
Vandermonde matrix ${\cal V}$.
%
\begin{gather}
  \uModal[] = \vand{}^{-1} \uNodal[]
\end{gather}

%%%%%%%%%%%%%%%%%%%%%%%%%%%%%%%%%%%%%%%%%%%%%%%%%%%%%%%%%%%%%%%%%%%%%%%%%%%%%%%%

After one (explicit Euler) timestep of size $\Delta t$ is applied to
\eref{EQN:1D_Advection}, the solution at the new time level is given by
\begin{gather}
  u^{\text{new}} (x)
= u              (x)
- \Delta t \frac{\partial}{\partial x} \Bigl[ f(x) \Bigr]
\end{gather}
In order for $u^\text{new}$ to be of the same polynomial order $\PO$ as $u$, 
the flux function $f$ must be one order higher.

%%%%%%%%%%%%%%%%%%%%%%%%%%%%%%%%%%%%%%%%%%%%%%%%%%%%%%%%%%%%%%%%%%%%%%%%%%%%%%%%

Since each $\uNodal[n]$ is simply a coefficient in \eref{EQN:Nodal_Form},
the derivative of this polynomial is just a linear combination of the
derivatives of the Lagrange polynomials.
%
\begin{gather}
  \label{EQN:Spatial_Derivative}
  \frac{d}{dr} u(r) = u_r(r)
                    = \sum_{n=1}^{N_P} \uNodal[n] \left[ \frac{d}{dr} l_n(r)
                                                  \right]
\end{gather}

%%%%%%%%%%%%%%%%%%%%%%%%%%%%%%%%%%%%%%%%%%%%%%%%%%%%%%%%%%%%%%%%%%%%%%%%%%%%%%%%

\subsubsection{Discontinuous Flux Function}%
\label{SEC:Intro:FR:DiscontinuousFlux}%
\key*[key:DiscontinuousFluxFunction]{}[discontinuous flux function]

%%%%%%%%%%%%%%%%%%%%%%%%%%%%%%%%%%%%%%%%%%%%%%%%%%%%%%%%%%%%%%%%%%%%%%%%%%%%%%%%

Using the nodal solution polynomial, we can evaluate the flux function at
each solution point, \thatis
$\fD[n] = \flxfun(\uNodal[n])$,
to obtain a polynomial of order $\PO$.
%
The superscript $D$ is used to identify it as the \emph{discontinuous} flux
function because the resulting function only uses information inside a given
cell and does not account for interaction with data from the neighboring cells.
%
Since we know the value of the discontinuous flux at each solution point in a
cell, we can compute its derivative using \eref{EQN:Spatial_Derivative}.
%
However, if we employ the derivative of $\fD{}$ to evaluate $f_x$, we obtain
erroneous solutions since such a derivative does not include the interaction
of data across cell interfaces.
%
Therefore, we need to construct a so-called \emph{continuous} flux function,
denoted by $\fC(r)$, which accounts for this interaction between adjacent cells
and approximates the discontinuous flux function in some sense,
to which we can then calculate its derivative.
%
The continuous flux function is obtained by adding a correction
to $\fD$, the discontinuous one.

%%%%%%%%%%%%%%%%%%%%%%%%%%%%%%%%%%%%%%%%%%%%%%%%%%%%%%%%%%%%%%%%%%%%%%%%%%%%%%%%

\subsubsection{Interface Common Flux Values}
\label{SEC:Intro:FR:Upwind}%
\key*[key:CommonFluxValue]{}[common flux value]


%%%%%%%%%%%%%%%%%%%%%%%%%%%%%%%%%%%%%%%%%%%%%%%%%%%%%%%%%%%%%%%%%%%%%%%%%%%%%%%%

One of the principles of a conservative scheme for steady-state flow
is that the flux leaving one cell should be the same as that entering
the next cell.
%
This is not the case for the discontinuous flux described above.
%
We therefore need to compute a common interface flux value that will be
used by the two adjacent cells.

%%%%%%%%%%%%%%%%%%%%%%%%%%%%%%%%%%%%%%%%%%%%%%%%%%%%%%%%%%%%%%%%%%%%%%%%%%%%%%%%

At each cell interface, let $\uLeft{}$ and $\uRight{}$
be the values taken on by the solution polynomials to its left and right,
respectively.
%
An upwind flux value that is common for the two adjacent cells, denoted by
$\fupw = \fupw (\uLeft,\uRight)$,
can then be defined via the flow direction.
%
Here, for the case of the Euler or Navier-Stokes equations, we use
Roe’s flux\cite{Roe_1981} with an entropy fix\cite{Huynh_1995}.

%%%%%%%%%%%%%%%%%%%%%%%%%%%%%%%%%%%%%%%%%%%%%%%%%%%%%%%%%%%%%%%%%%%%%%%%%%%%%%%%

\subsubsection{Left\slash Right Flux Jumps}
\label{SEC:Intro:FR:FluxJumps}

%%%%%%%%%%%%%%%%%%%%%%%%%%%%%%%%%%%%%%%%%%%%%%%%%%%%%%%%%%%%%%%%%%%%%%%%%%%%%%%%

Let \fupwL{} denote the common upwind flux at the `left boundary' of the cell
and \fupwR{} denote the common upwind flux at the `right boundary' of the cell.
%
There is a discrepancy between the value of the discontinuous flux function
evaluated at the left and right boundaries of the cell and the desired
common upwind flux values.
%
We denote these left\slash right flux jumps as
$\fluxjumpL{} = \fupwL{} - \fD\pvp[-1]$ and
$\fluxjumpR{} = \fupwR{} - \fD\pvp[ 1]$.

 
%%%%%%%%%%%%%%%%%%%%%%%%%%%%%%%%%%%%%%%%%%%%%%%%%%%%%%%%%%%%%%%%%%%%%%%%%%%%%%%%

\subsubsection{Correction Functions}
\label{SEC:Intro:FR:Correction_Functions}%
\key*[key:CorrectionFunction]{}[correction function]

%%%%%%%%%%%%%%%%%%%%%%%%%%%%%%%%%%%%%%%%%%%%%%%%%%%%%%%%%%%%%%%%%%%%%%%%%%%%%%%%

Our goal is to construct a continuous flux function, $\fC(r)$, as a polynomial
of degree $\PO+1$ by requiring that it takes on the common upwind flux values
at the two cell interfaces, and it approximates $\fD$ within the interior of
the cell via the $\PO-1$ remaining constraint conditions for the polynomial.
%
However, instead of defining \fC{} directly, we define $\G=\fC-\fD$, which
represents the \emph{corrected flux function} that must be added to
the \key[key:DiscontinuousFluxFunction]{discontinuous flux function}[=] $\fD$
in order to obtain
the \key[key:ContinuousFluxFunction]{continuous flux function}[=] $\fC$.

 
At the two interfaces, $r=\pm{}1$, the function $\G$ takes the following
values which are the flux jumps described previously.
%
\begin{gather}
  \G\pvp[-1] = \fC\pvp[-1] - \fD\pvp[-1]
             = \fupwL      - \fD\pvp[-1] = \fluxjumpL{} \\
  \G\pvp[+1] = \fC\pvp[+1] - \fD\pvp[+1]
             = \fupwR      - \fD\pvp[+1] = \fluxjumpR{}
\end{gather}
%
Therefore, $\G\pvp[r]$ has prescribed left and right interface values,
is of degree $\PO+1$, and approximates zero within the interior of the cell.

%%%%%%%%%%%%%%%%%%%%%%%%%%%%%%%%%%%%%%%%%%%%%%%%%%%%%%%%%%%%%%%%%%%%%%%%%%%%%%%%

There are many ways one could prescribe how the correction approximates
the zero function, and different choices result in different
numerical schemes with different stability characteristics.
%
Further discussion regarding the choices in correction functions can be found
in \linecite{Huynh_2007,Vincent_2011,Castonguay_2012,Vincent_2015}.
%
The GFR code uses \key{Radau polynomials}[polynomial!Radau]
as the correction functions.
%
This choice has been shown to recover a nodal discontinuous Galerkin (DG)
method.

%%%%%%%%%%%%%%%%%%%%%%%%%%%%%%%%%%%%%%%%%%%%%%%%%%%%%%%%%%%%%%%%%%%%%%%%%%%%%%%%

We separate the prescription of the jump at the left interface from that of
the right.
%
This separation plays a critical role in the approach.
%
\begin{align}
   \G\pvp[r] =& \GL\pvp[r]   +  \GR\pvp[r] &
  \GL\pvp[r] =& \fluxjumpL{} \, \gL\pvp[r] &
  \GR\pvp[r] =& \fluxjumpR{} \, \gR\pvp[r]
\end{align}
%
Let \gL{} be the \key{\emph{correction function}}[correction function]
for the left interface of the reference element defined by
%
\begin{equation}
  \gL\pvp[-1] = 1, \qquad
  \gL\pvp[ 1] = 0
\end{equation}
%
It is a polynomial of degree $\PO+1$ approximating the zero function in
some sense.
%
We always choose \gR{} by reflection: $\gR\pvp[\rv]=\gL\pvp[-\rv]$.
%
As a result, it is also a polynomial of degree $\PO+1$ with
boundary values of
%
\begin{equation}
  \gR\pvp[-1] = 0, \qquad
  \gR\pvp[ 1] = 1
\end{equation}

%%%%%%%%%%%%%%%%%%%%%%%%%%%%%%%%%%%%%%%%%%%%%%%%%%%%%%%%%%%%%%%%%%%%%%%%%%%%%%%%

In GFR, the left\slash right correction function of order $\PO+1$ is computed
using the right\slash left Radau polynomial of that same order.
%
Details of the Radau polynomial may be found in \aref{SEC:RadauPoly}.
%
\Fref{FIG:RadauPoly} illustrates the shape of different orders of the
Radau polynomial.
%
The value at one side is always zero, and the value at the other side
is always one.
%
The number of times the polynomial equals zero on \stdelem{} is equal to the
order of the polynomial.
%
This choice results in a correction that approximates the zero function
in a least-squares sense, and the resulting method is equivalent to a
nodal discontinuous Galerkin (DG) scheme.

%%%%%%%%%%%%%%%%%%%%%%%%%%%%%%%%%%%%%%%%%%%%%%%%%%%%%%%%%%%%%%%%%%%%%%%%%%%%%%%%

\begin{figure}
  \begin{centering}
    \begin{subfigure}{3.0in}
      \insertimage[width=\textwidth]{Intro/RadauLeft}
      \caption{Left Radau Polynomial}
      \label{FIG:LeftRadauPoly}
    \end{subfigure}%
    \begin{subfigure}{3.0in}
      \insertimage[width=\textwidth]{Intro/RadauRight}
      \caption{Right Radau Polynomial}
      \label{FIG:RightRadauPoly}
    \end{subfigure}
    \caption{Left and Right \key{Radau polynomials}[polynomial!Radau].}
    \label{FIG:RadauPoly}
  \end{centering}
\end{figure}

%%%%%%%%%%%%%%%%%%%%%%%%%%%%%%%%%%%%%%%%%%%%%%%%%%%%%%%%%%%%%%%%%%%%%%%%%%%%%%%%

%%%%%%%%%%%%%%%%%%%%%%%%%%%%%%%%%%%%%%%%%%%%%%%%%%%%%%%%%%%%%%%%%%%%%%%%%%%%%%%%

\subsubsection{Continuous Flux Function}
\label{SEC:Intro:FR:ContinuousFlux}%
\key*[key:ContinuousFluxFunction]{}[continuous flux function]


The \emph{continuous flux function} may now be computed by adding
the \key[key:CorrectionFunction]{correction function}[=] to
the \key[key:DiscontinuousFluxFunction]{discontinuous flux function}[=].
\begin{align}
  \fC =& \fD + \G \\
      =& \fD + \GL + \GR \\
      =& \fD + \fluxjumpL{} \, \gL + \fluxjumpR{} \, \gR
\end{align}

As will become apparent in the following two steps,
it is not necessary to explicitly compute the values of the
continuous flux function.
%
We need only know how it is composed so that we can compute its derivative.

%%%%%%%%%%%%%%%%%%%%%%%%%%%%%%%%%%%%%%%%%%%%%%%%%%%%%%%%%%%%%%%%%%%%%%%%%%%%%%%%

\subsubsection{Derivative of the Continuous Flux Function}%
\label{SEC:Intro:FR:DerivativeFlux}
\key*[key:ContinuousFluxFunctionDerivative]{}%&
     [continuous flux function!derivative]

The derivative of the continuous flux function is given by:
%
\begin{align}
     \frac{\partial \fC}{\partial r}
  =& \frac{\partial \fD}{\partial r}
  +  \fluxjumpL{} \, \gL' + \fluxjumpR{} \, \gR'
\end{align}
%
where $\partial \fD/\partial r$ can be computed in a manner similar to
\eref{EQN:Spatial_Derivative} and expressions for $\gL'$ and $\gR'$
follow from the derivatives of the Radau polynomials found in
\aref{SEC:RadauPoly}.
%
The inverse mapping function $\chi^{-1}$ can then be used to transform
the derivative from reference space ($r$) to physical space ($x$),
and that value for $\partial F/\partial x$ can be used to approximate
$\partial f/\partial x$ at each solution point within the grid cell.

%%%%%%%%%%%%%%%%%%%%%%%%%%%%%%%%%%%%%%%%%%%%%%%%%%%%%%%%%%%%%%%%%%%%%%%%%%%%%%%%

\subsubsection{Time Advancement}
\label{SEC:Intro:FR:Time_Advancement}

%%%%%%%%%%%%%%%%%%%%%%%%%%%%%%%%%%%%%%%%%%%%%%%%%%%%%%%%%%%%%%%%%%%%%%%%%%%%%%%%

After one (explicit Euler) timestep of size $\Delta t$ is applied to
\eref{EQN:1D_Advection}, the solution at each point $n$ at the new time level
is given by
\begin{gather}
  u_n^{\text{new}} (x)
= u_n              (x)
- \Delta t \left. \frac{\partial}{\partial x} \Bigl[ F(x) \Bigr]
           \right|_n
\end{gather}
%
More advanced time stepping algorithms work similarly.
%
Further details on time-stepping in GFR can be found in \sref{SEC:Time}.

%%%%%%%%%%%%%%%%%%%%%%%%%%%%%%%%%%%%%%%%%%%%%%%%%%%%%%%%%%%%%%%%%%%%%%%%%%%%%%%%

\subsection{Solution Points}
\label{SEC:Intro:Solution_Points}
\key*{}[points!solution]

%%%%%%%%%%%%%%%%%%%%%%%%%%%%%%%%%%%%%%%%%%%%%%%%%%%%%%%%%%%%%%%%%%%%%%%%%%%%%%%%

As discussed in \sref{SEC:Intro:FR:Polynomial}, the solution within each
element is represented by a polynomial of order $\PO$ defined by values
at $\PO[+1]$ solution points.
%
\key{Equally-spaced (equidistant)}[points!equidistant],
\key{Lobatto}[points!Legendre-Gauss-Lobatto (Lobatto)], or
\key{Gauss}[points!Legendre-Gauss (Gauss)] points may be used.
%
\Fref{FIG:1D_Points} illustrates the distribution of these points
within a standard element in one dimension.
%
In each case, the distribution of points within the reference element
is symmetric about the center of the cell.
%
Equidistant points include points on both boundaries of the cell
and uniform spacing is used for the interior points.
%
Lobatto points also include points on both boundaries of the cell,
but the interior spacing is non-uniform and biased toward the
boundaries of the cell.
%
Gauss points are distributed purely within the interior of the cell,
and are also biased toward both boundaries of the cell.

%%%%%%%%%%%%%%%%%%%%%%%%%%%%%%%%%%%%%%%%%%%%%%%%%%%%%%%%%%%%%%%%%%%%%%%%%%%%%%%%

\begin{figure}
  \begin{centering}
    \begin{subfigure}{2.0in}
      \insertimage[width=\textwidth]{Intro/1D_Points_Equal}
      \caption{\key*{Equidistant points}[points!equidistant].}
      \label{FIG:1D_Points_Equal}
    \end{subfigure}%
    \begin{subfigure}{2.0in}
      \insertimage[width=\textwidth]{Intro/1D_Points_Lobatto}
      \caption{\key*{Lobatto points}[points!Legendre-Gauss-Lobatto (Lobatto)].}
      \label{FIG:1D_Points_Lobatto}
    \end{subfigure}%
    \begin{subfigure}{2.0in}
      \insertimage[width=\textwidth]{Intro/1D_Points_Gauss}
      \caption{\key*{Gauss points}[points!Legendre-Gauss (Gauss)].}
      \label{FIG:1D_Points_Gauss}
    \end{subfigure}
    \caption{1D solution point distributions vs polynomial order,
             $\PO[1]-\PO[4]$.}
    \label{FIG:1D_Points}
  \end{centering}
\end{figure}

%%%%%%%%%%%%%%%%%%%%%%%%%%%%%%%%%%%%%%%%%%%%%%%%%%%%%%%%%%%%%%%%%%%%%%%%%%%%%%%%

\begin{figure}
  \begin{centering}
    \begin{subfigure}{1.5in}
      \insertimage[width=\textwidth]{Triangle_Pts/Triangle_Gauss_P1}
      \caption{\PO[1]}
      \label{FIG:Triangle_Gauss_P1}
    \end{subfigure}%
    \begin{subfigure}{1.5in}
      \insertimage[width=\textwidth]{Triangle_Pts/Triangle_Gauss_P2}
      \caption{\PO[2]}
      \label{FIG:Triangle_Gauss_P2}
    \end{subfigure}%
    \begin{subfigure}{1.5in}
      \insertimage[width=\textwidth]{Triangle_Pts/Triangle_Gauss_P3}
      \caption{\PO[3]}
      \label{FIG:Triangle_Gauss_P3}
    \end{subfigure}
    \begin{subfigure}{1.5in}
      \insertimage[width=\textwidth]{Triangle_Pts/Triangle_Gauss_P4}
      \caption{\PO[4]}
      \label{FIG:Triangle_Gauss_P4}
    \end{subfigure}
    \caption{2D triangle \quoterm{Gauss} solution point distributions vs
             polynomial order, $\PO[1]-\PO[4]$.}
    \label{FIG:2D_Triangle_GaussPoints}
  \end{centering}
\end{figure}

%%%%%%%%%%%%%%%%%%%%%%%%%%%%%%%%%%%%%%%%%%%%%%%%%%%%%%%%%%%%%%%%%%%%%%%%%%%%%%%%

Gauss points are generally preferred over Lobatto because the numerical
quadrature associated with Gauss points has a slightly higher degree of
accuracy than that for Lobatto points.
%
Equidistant points should not be used in the FR method due to quadrature
stability issues; they should only be used for post-processing\slash
visualization purposes.
%
\Fref{FIG:2D_Triangle_GaussPoints} illustrates the distribution of
\quoterm{Gauss} points within a 2D triangle element.
%
A tensor product is used to distribute solution points within 2D quadrilaterals,
which is simply a 1D distribution in each direction.
%
Similarly, a 3D prism and 3D hexahedron use a 1D distribution in the
third direction.
%
The distribution in a 3D tetrahedron and 3D pyramid are somewhat similar
to the 2D triangle, but are difficult to visualize on a printed page.

%%%%%%%%%%%%%%%%%%%%%%%%%%%%%%%%%%%%%%%%%%%%%%%%%%%%%%%%%%%%%%%%%%%%%%%%%%%%%%%%

\SeeAlso{
\begin{tabular}[t]{lll}
  \inp[nmlNumerics::AllOrders]{All\_Orders}
& \inp[nmlNumerics::AllPoints]{All\_Points}
& \inp[nmlNumerics::LocTrianglePts]{Loc\_Triangle\_Pts}
\\
  \inp[nmlNumerics::SolutionOrder]{Solution\_Order}
& \inp[nmlNumerics::LocSolutionPts]{Loc\_Solution\_Pts}
& \inp[nmlNumerics::LocTetrahedraPts]{Loc\_Tetrahedra\_Pts}
\\
  \inp[nmlNumerics::ProjectionOrder]{Projection\_Order}
& \inp[nmlNumerics::LocFluxPts]{Loc\_Flux\_Pts}
\\
  \inp[nmlNumerics::QuadratureOrder]{Quadrature\_Order}
& \inp[nmlNumerics::LocQuadraturePts]{Loc\_Quadrature\_Pts}
\\
  \inp[nmlNumerics::ErrorOrder]{Error\_Order}
& \inp[nmlNumerics::LocErrorPts]{Loc\_Error\_Pts}
%
\end{tabular}
}

%%%%%%%%%%%%%%%%%%%%%%%%%%%%%%%%%%%%%%%%%%%%%%%%%%%%%%%%%%%%%%%%%%%%%%%%%%%%%%%%

\subsection{Quadrature Points}
\label{SEC:Intro:Quadrature_Points}
\key*{}[points!quadrature]

%%%%%%%%%%%%%%%%%%%%%%%%%%%%%%%%%%%%%%%%%%%%%%%%%%%%%%%%%%%%%%%%%%%%%%%%%%%%%%%%

Instead of directly evaluating the nonlinear flux at the solution points,
it can be more accurate to evaluate the flux using a higher quadrature order
(\thatis a higher order polynomial and its corresponding points)
and then projecting that onto a polynomial of the same order as the solution.
%
\Sref{SEC:Aliasing} provides a more detailed explanation of this
aliasing error and the steps involved in using quadrature points.

%%%%%%%%%%%%%%%%%%%%%%%%%%%%%%%%%%%%%%%%%%%%%%%%%%%%%%%%%%%%%%%%%%%%%%%%%%%%%%%%

\subsection{Output Points}
\label{SEC:Intro:Output_Points}
\key*{}[points!output]

%%%%%%%%%%%%%%%%%%%%%%%%%%%%%%%%%%%%%%%%%%%%%%%%%%%%%%%%%%%%%%%%%%%%%%%%%%%%%%%%

GFR allows the user to output the solution at a different set of points
(i.e., \key{equidistant}[points!equidistant]) for visualization.
%
The reason for this is that since
\key{Gauss points}[points!Legendre-Gauss (Gauss)] do not extend to the
boundaries of cells, many visualization tools will render the result
with a small gap between each cell.

\IDEA{
Include 2D contour plot using output points at Gauss points and equidistant.
}

\SeeAlso{
\inp[nmlNumerics::AllOrders]{All\_Orders},
\inp[nmlOutput::OutputOrder]{Output\_Order},
\inp[nmlOutput::LocOutputPts]{Loc\_Output\_Pts}
}

%%%%%%%%%%%%%%%%%%%%%%%%%%%%%%%%%%%%%%%%%%%%%%%%%%%%%%%%%%%%%%%%%%%%%%%%%%%%%%%%

%\subsection{H-Refinement vs \PO-Refinement}
\subsection{Solution Accuracy and Grid Refinement}
\label{SEC:Intro:HP_Refinement}

%%%%%%%%%%%%%%%%%%%%%%%%%%%%%%%%%%%%%%%%%%%%%%%%%%%%%%%%%%%%%%%%%%%%%%%%%%%%%%%%

%In discussing H- versus \PO-refinement, one needs an objective measure of
%computational cost.
%
We start by introducing the term
\key*{degree-of-freedom (DoF)}[degree-of-freedom (DoF)] which is simply a
location at which we are solving the governing equations.
%
Note, some refer to a DoF as an individual equation being solved at a point, so
an individual solution point would contain 5 DoF for the 3D Euler\slash
Navier-Stokes equations; here we refer to a DoF as the solution point itself.
%
For traditional finite-difference and finite-volume methods, the governing
equations are solved at each grid node (or cell for cell-based methods), so the
total DoF is just the total number of nodes (or cells) in the grid.
%
Within the context of high-order finite-\slash spectral-element methods, a
single grid cell will solve the governing equations on a set of solution nodes
(or modes for modal methods) which we refer to as solution points.
%
Each of these solution points is a DoF in a simulation, and the total number of
solution points in a cell depends on the element type and the polynomial order
of the solution (\sref{SEC:Element_Types} provides more detail on the DoF for each
element type).
%
Assuming ideal parallel scaling with minimal communication costs, the
computational cost of a simulation is proportional to the total DoF within the
entire computational domain.
%
With finite-difference, finite-volume, and spectral-element methods being so
vastly different, the fairest way to compare the computational costs of the
different methods is only between simulations containing similar numbers of
total DoF.

%\FIXME{
%\begin{enumerate}
%\item
%Is DoF really an accurate measure of total computational cost,
%or only a measure of (ie, matrix multiply) operation count?
%\item
%Does DoF account for reduced operation count of sparse quad\slash hex
%matrices?
%\item
%How do h\slash p effects on timestep effect cost, and do they scale differently?
%\end{enumerate}
%}

%%%%%%%%%%%%%%%%%%%%%%%%%%%%%%%%%%%%%%%%%%%%%%%%%%%%%%%%%%%%%%%%%%%%%%%%%%%%%%%%

%H-refinement refers to modifying the physical mesh to better
%resolve flow features.
%%
%For a uniform 1D mesh, this spacing is often denoted as $h=\Delta x$,
%and thus the name.
%%
%For complex 3D cases with nonuniform element types \andor grid spacing,
%this refinement is often done by halving the edge length of each
%element and creating additional sub-elements.
%
%In high-order CFD, \PO-refinement refers to increasing the order $\PO$
%of the solution polynomial used within each element, without changing
%the physical mesh.
%%
%In many cases, \PO-refinement can achieve higher accuracy with fewer
%degrees of freedom and therefore at a lower computational cost.
%%
%Another benefit is that it does not require any remeshing.

%%%%%%%%%%%%%%%%%%%%%%%%%%%%%%%%%%%%%%%%%%%%%%%%%%%%%%%%%%%%%%%%%%%%%%%%%%%%%%%%

The computational mesh plays a critical role in any scale-resolving simulation
(SRS), as the grid will filter out the scales of turbulent structures smaller
than the grid resolution.
%
As such, it is important to understand the effect that refinement---grid ($h$),
order ($p$), or combination thereof ($h/p$)---will have on any given simulation.
%
For finite-difference and finite-volume methods, $h$-refinement is only possible
by changing the number of grid nodes\slash cells.
%
True $p$-refinement is only possible by increasing the order of accuracy while
keeping the grid resolution constant, \forexample comparing \nth{2}- and
\nth{4}-order finite-difference schemes on the same structured grid.
%
For spectral-element methods like FR, increasing $\PO$ on the same grid
results in $h/p$-refinement because the number of DoF in each cell also
increases with \PO.
%
Therefore, in the context of FR, $p$-refinement requires decreasing the grid
resolution such that the total DoF across all grid cells is equivalent between
different values of \PO.

%%%%%%%%%%%%%%%%%%%%%%%%%%%%%%%%%%%%%%%%%%%%%%%%%%%%%%%%%%%%%%%%%%%%%%%%%%%%%%%%

Achieving quality wall-resolved LES (WRLES) results can be a tricky and time
consuming process that is highly dependent on the grid quality and numerical set
up.
%
WRLES generally requires high-resolution meshes for both accuracy and stability,
so much so that the maximum stable time step can be multiple orders of magnitude
smaller compared to RANS.
%
One of the major benefits of spectral-element methods is that the order of the
method can be increased while simultaneously decreasing the size of the grid to
maintain a certain level of accuracy of a simulation.
%
Decreasing the size of the grid will generally increase the size of the smallest
grid cells which will directly increase the maximum allowable stable time step;
however, increasing \PO{} also reduces the maximum allowable stable time step
since the time step scales as $1/\PO^2$. 
%
Furthermore, increasing \PO{} can introduce numerical instabilities that require
added numerical dissipation to prevent the simulation from diverging.
%
Consequently, this numerical dissipation will also damp out the smallest
turbulent scales within a simulation unless the dissipation can be precisely
targeted to only the areas within the solution where it is absolutely needed.
%
Ultimately, there is a back and forth required with grid and numerical setups to
find an optimal combination to achieve accuracy goals within the computational
time available for a set of simulations.


%%%%%%%%%%%%%%%%%%%%%%%%%%%%%%%%%%%%%%%%%%%%%%%%%%%%%%%%%%%%%%%%%%%%%%%%%%%%%%%%

